\section{Judgment kept apart from money}
\label{sec:27.5}
The third piece is who judges. Somebody has to decide which protocols are sound, which allocations go ahead and which awards are made, and the status body has to decide what a kind of system's architecture shows. Those are qualitative judgments, and the book's rule for them is the one it uses for the adverse quorum: panels drawn by lot from pools the Association's members nominate, with practitioners, members of the classes the charter names, and scholars among them. The arts endowment's panels are the nearest precedent. By statute they are to reflect “a wide geographic, ethnic, and minority representation as well as individuals reflecting diverse artistic and cultural points of view,” to include lay members who know the arts and don't practice them as a profession, and to change substantially from year to year, with nobody serving more than three years running \autocite{us1965artsact}. Those panels are chosen under the chair's procedures. Here the lot chooses.

What the panels judge is kinds, methods and protocols. They never judge a particular student's mind: a disposition is presumed the student's own, and only the record of an installation rebuts that. The bodies stay apart. The status body finds what is true of kinds of system; the review board decides what may be done to a candidate to find it out; the panels allocate and award; the procurement body and the Treasury pay. None decides another's question, and the bodies that pay never judge.

The method by which a polity learns what is true, about these systems as about anything, is the scientific one, and it is a social method. Longino's account of objectivity is the one this design follows: objectivity belongs to a community's practices of criticism and not to any one inquirer, and it depends on venues where criticism can be made, standards the critics share, a community that takes criticism up, and intellectual authority shared among those qualified to give it \autocite{longino1990science}. That is the reason the people these systems act on sit on the panels, and the reason the lot picks them. A community of inquiry that keeps its critics out has lost what made it objective. It is also Dewey's point about the public, and Anderson's about democracy as a way of inquiring (the next chapter): the people who bear a problem know things about it that nobody else does \autocite{dewey1927public,anderson2006epistemology}.
